\section{Modelos de comunicación}

\begin{frame}{Spotify: comunicación cliente-servidor}
  \begin{itemize}
    \item \textbf{Cliente--servidor:} las aplicaciones envían solicitudes a
    los servicios de Spotify. El perímetro recibe el tráfico y lo dirige al
    backend \cite{spotifyIncident2025}.
    \item \textbf{Solicitud--respuesta:} en la Web API pública, un cliente
    autorizado envía \texttt{GET /v1/search} mediante HTTPS y recibe los
    resultados del catálogo en JSON. La respuesta corresponde a esa solicitud \cite{spotifyWebApiSearch}.
    \item \textbf{Mensajería asíncrona:} la plataforma de datos recopila
    eventos de uso. Google Cloud Pub/Sub desacopla la publicación de eventos
    de su consumo y procesamiento posterior \cite{khlebnikova2024dataplatform}.
  \end{itemize}
\end{frame}

\begin{frame}{Flujo de solicitudes y eventos}
  \centering
  \IfFileExists{figuras/flujo_comunicacion.png}{%
    \includegraphics[width=\linewidth,height=0.70\textheight,keepaspectratio]{flujo_comunicacion.png}%
  }{%
    \begin{itemize}
      \item \textbf{Consulta:} el cliente envía una solicitud HTTPS a la
      Web API de Spotify y recibe una respuesta JSON.
      \item \textbf{Eventos:} el cliente emite un evento de uso. El servicio
      de eventos lo publica en Pub/Sub y los consumidores lo procesan después.
    \end{itemize}%
  }%
  \par\smallskip
  {\scriptsize Elaboración propia a partir de Spotify for Developers y
  Spotify Engineering \cite{spotifyWebApiSearch,khlebnikova2024dataplatform}.\par}
\end{frame}
